%% evaluation.tex
%%

%% ==================
\section{Evaluation}
\label{sec:Evaluation}
%% ==================

Table~\ref{tab:results} reports mean $\pm$ standard deviation across 5 seeds (42-46) on a pooled, held-out test set drawn from all 6 sites.

\begin{table}[htp]
\centering
\caption{Results across 5 seeds (mean $\pm$ std)}
\label{tab:results}
\begin{tabular}{lccc}
\hline
Model & Accuracy & Macro-F1 & Critical Recall (class 2) \\
\hline
Centralized RF (baseline) & 0.9995 $\pm$ 0.0005 & 0.9996 $\pm$ 0.0005 & 0.9994 $\pm$ 0.0013 \\
Federated Ensemble RF (improved) & 0.9959 $\pm$ 0.0024 & 0.9949 $\pm$ 0.0024 & 0.9997 $\pm$ 0.0007 \\
Single-Site Local Only (no federation) & 0.9761 $\pm$ 0.0046 & 0.9737 $\pm$ 0.0051 & 0.9958 $\pm$ 0.0031 \\
\hline
\end{tabular}
\end{table}

%% ===============================
\subsection{Experiment 1: The Cost of Going It Alone}
\label{sec:evalOne}
%% ===============================

A single site's own local model, used in isolation with no federation and no centralization, reaches only 97.4\% macro-F1 -- a clear, consistent gap below both the centralized (99.96\%) and federated (99.49\%) variants. This confirms the premise of the whole exercise: because each site's local traffic mixture is skewed toward one of the three scenarios (uniform/burst/real-time), a site that only ever sees its own data generalises worse to the full population it will actually encounter in production.

%% ===============================
\subsection{Experiment 2: Federation Recovers Most of the Centralization Benefit}
\label{sec:evalTwo}
%% ===============================

The federated ensemble closes the large majority of the gap between local-only and centralized (macro-F1 97.37\% $\to$ 99.49\%, versus centralized's 99.96\%), while never pooling raw telemetry across sites -- directly addressing the privacy/centralization limitation the baseline paper names in its own Discussion section. The remaining gap to full centralization (99.49\% vs. 99.96\% macro-F1, about half a point) is the measurable cost of keeping data decentralized, and is small relative to the 2.6-point gap federation recovers from the single-site baseline.

%% ===============================
\subsection{Experiment 3: The Minority Class}
\label{sec:evalThree}
%% ===============================

Critical Recall (the class the baseline paper itself needed SMOTE to handle, at only 9.64\% of their real-world data) is high and statistically indistinguishable across all three variants (99.6-100\%), including the single-site baseline. This suggests SMOTE-balancing within each site is sufficient on its own to protect the minority class here; federation's main benefit in this experiment is on overall accuracy/macro-F1 (driven by better generalisation across traffic regimes), not on minority-class recall specifically.

%% ===============================
\subsection{Discussion}
\label{sec:discussion}
%% ===============================

All three variants score unusually high in absolute terms (97-100\%) compared to the baseline paper's own reported 96-98\% for its best models. This is expected: the synthetic generator's labelling rule is a deterministic function of the same features the models are trained on, with less real-world sensor noise than Azure IoT telemetry would have. The relative ordering of the three variants -- local-only worst, federated in the middle, centralized best -- is the finding that matters here, not the absolute numbers, and that ordering is what a privacy/utility trade-off in federated learning would predict.
